%%%%%%%%%%%%%%%%%%%%%%%%%%%%%%%%%%%%%%%%%%%%%%%%%%%%%%%%%%%%%%%%%%%%%%
%
% Discrete Appell-Dunkl sequences 
% and Bernoulli-Dunkl polynomials of the second kind
% by
% J. I. Extremiana, E. Labarga, J. M\'inguez and J. L. Varona
%
% email: jextremi@unirioja.es, edgar.labarga@unirioja.es,
%        judit.minguez@unirioja.es, jvarona@unirioja.es
%
% COMIENZO DE ESCRITURA: 2021/05/05
% LAST MODIFICATION: 2021/05/19
%
%%%%%%%%%%%%%%%%%%%%%%%%%%%%%%%%%%%%%%%%%%%%%%%%%%%%%%%%%%%%%%%%%%%%%%

%% ORCID:
% J. I. Extremiana: 0000-0002-6224-7768
% E. Labarga: 0000-0003-1017-869X
% J. M\'inguez: 0000-0002-9840-5772
% J. L. Varona: 0000-0002-2023-9946

\documentclass[11pt]{article}
%\documentclass[10pt,a4paper]{article}

\usepackage{amsmath, amsthm, amsfonts}

\addtolength{\textwidth}{1.6cm}
\addtolength{\oddsidemargin}{-0.8cm}
\setlength{\evensidemargin}{\oddsidemargin}

%\usepackage{graphicx}

\usepackage{hyperref}

%%% ONLY FOR DRAFTS:
%\usepackage[utf8]{inputenc}
%\usepackage[notcite,notref]{showkeys}
%\usepackage{color}
%\newcommand{\rojo}[1]{\textcolor{red}{#1}}
%\newcommand{\azul}[1]{\textcolor{blue}{#1}}

%%% ONLY FOR DRAFTS (put it after hyperref)
%\usepackage{backref}
%\usepackage{refcheck}


%%% \allowdisplaybreaks es una orden de amsmath que permite cortes de
%%% p\'agina dentro de una f\'ormula de varias l\'ineas (align, multline, gather, ...)
\allowdisplaybreaks[3]
%%% Escala de 1 (permite algo) a 4 (ninguna traba para cortar).
%%% Si en alg\'un sitio queremos impedir un posible corte, usar \\* en vez de \\

\theoremstyle{plain}
\newtheorem{theorem}{Theorem}[section]
\newtheorem{proposition}[theorem]{Proposition}
\newtheorem{lemma}[theorem]{Lemma}
\newtheorem{corollary}[theorem]{Corollary}

\theoremstyle{remark}
\newtheorem*{remark}{Remark}

\numberwithin{equation}{section} % To include the number of section

\newcommand{\Ea}{E_{\alpha}}
\newcommand{\La}{\Lambda_{\alpha}}
\newcommand{\Lax}{\Lambda_{\alpha,x}}
\newcommand{\Fa}{\mathcal{F}_{\alpha}} % Dunkl transform of order \alpha
\newcommand{\mua}{\mu_{\alpha}}
\newcommand{\binoma}[2]{\binom{#1}{#2}_{\!\alpha}}
\newcommand{\arcsinh}{\operatorname{arcsinh}}

\newcommand{\I}{\mathcal{I}}

\newcommand{\B}{\mathfrak{B}}
%\newcommand{\E}{\mathcal{E}}

\newcommand{\NN}{\mathbb{N}}
\newcommand{\RR}{\mathbb{R}}
\newcommand{\ZZ}{\mathbb{Z}}
\newcommand{\CC}{\mathbb{C}}
\newcommand{\ZZO}{\mathbb{Z} \setminus \{0\}}

%\renewcommand{\Im}{\operatorname{Im}}
%\renewcommand{\Re}{\operatorname{Re}}
%\newcommand{\sgn}{\operatorname{sgn}}

\newcommand{\cc}{\gamma}

%----------
\title{Discrete Appell-Dunkl sequences 
and Bernoulli-Dunkl polynomials of the second kind%
\footnote{\textbf{This paper has been published as:}
J. I. Extremiana Aldana, E. Labarga, J. M\'inguez Ceniceros and J. L. Varona, 
Discrete Appell-Dunkl sequences and Bernoulli-Dunkl polynomials of the second kind,
\textit{J.\ Math.\ Anal.\ Appl.}\ \textbf{507} (2022), no.~2, 125832, 20 pp., 
\url{https://doi.org/10.1016/j.jmaa.2021.125832}}
\footnote{Partially supported by grant PGC2018-096504-B-C32 
(Minis\-te\-rio de Eco\-no\-m\'ia y Com\-pe\-ti\-ti\-vi\-dad, Spain)}}
%----------


%----------
\author{Jos\'e Ignacio Extremiana Aldana\textsuperscript{1}, 
Edgar Labarga\textsuperscript{2}, 
Judit M\'inguez Ceniceros\textsuperscript{3},\\
and Juan Luis Varona\textsuperscript{4}\\[6pt]
\small Departamento de Matem\'aticas y
Computaci\'on, Universidad de La Rioja,\\[-2pt]
\small 26006 Logro\~no, Spain\\[3pt]
\small \textsuperscript{1}{jextremi@unirioja.es},
\textsuperscript{2}{edgar.labarga@unirioja.es},\\
\small \textsuperscript{3}{judit.minguez@unirioja.es},
\textsuperscript{4}{jvarona@unirioja.es}}
%----------

\date{}

%----------
\begin{document}
%----------

%----------
\maketitle
%----------

\begin{abstract}
In a similar way that the Appell sequences of polynomials can be extended to the Dunkl context,
where the ordinary derivative is replaced by Dunkl operator on the real line, and the exponential function
is replaced by the so-called Dunkl kernel, one can expect that the discrete Appell sequences can be
extended to the Dunkl context. In this extension, the role of the ordinary translation is played by
the Dunkl translation, that is a much more intricate operator.
 
In this paper, we define discrete Appell-Dunkl sequences of polynomials, and we give some properties 
and examples. In particular, we show which is the suitable definition for the Bernoulli polynomials 
of the second kind in the Dunkl context.

2010 Mathematics Subject Classification: Primary 11B68; Secondary 33C10, 11B83.

Keywords: Appell-Dunkl sequences, Bernoulli-Dunkl polynomials, Bernoulli-Dunkl polynomials of the second kind, Dunkl transform.
\end{abstract}

%----------
\section{Introduction}
%----------

An Appell sequence $\{P_k(x)\}_{k=0}^{\infty}$ is a sequence of polynomials such that 
\begin{equation}
\label{eq:A}
  \frac{d}{dx} P_k(x) = k P_{k-1}(x), \quad k \ge 1.
\end{equation}
If instead of the derivative we use the discrete operator $\Delta f(x) = f(x+1) - f(x)$, we say that a discrete Appell sequence $\{p_k(x)\}_{k=0}^{\infty}$ is a sequence of polynomials such that 
\begin{equation}
\label{eq:Adiscrete}
  p_k(x+1) - p_k(x) = k p_{k-1}(x), \quad k \ge 1.
\end{equation}
It is well known that Appell sequences can be defined by a Taylor generating expansion
\begin{equation}
\label{eq:As}
  A(t) e^{xt} = \sum_{k=0}^{\infty} P_k(x) \frac{t^k}{k!},
\end{equation}
where $A(t)$ is a function analytic at $t=0$ with $A(0) \ne 0$; 
similarly, discrete Appell sequences can be defined by a Taylor generating expansion
\begin{equation}
\label{eq:dAs}
  A(t)(1+t)^x = \sum_{k=0}^{\infty} p_k(x) \frac{t^k}{k!},
\end{equation}
where $A(t)$ is a function analytic at $t=0$ with $A(0) \ne 0$. 

There is a wide mathematical literature studying families of Appell sequences.
Typical examples are the trivial case $\{x^k\}_{k=0}^{\infty}$ 
whose generating function is \eqref{eq:As} with $A(t)=1$,
as well as the Bernoulli polynomials that were used by Euler in 1740 to sum 
$\sum_{n=1}^{\infty} 1/n^{2k}$, and whose generating function is
\eqref{eq:As} with $A(t)=t/(e^t-1)$.

In the case of discrete Appell sequences, the trivial case,
obtained from \eqref{eq:dAs} with $A(t)=1$,
is the family $\{x^{\underline{k}}\}_{k=0}^{\infty}$ where
\begin{equation}
\label{eq:ffact}
x^{\underline{k}} = x(x-1)  \cdots (x-k+1) = \prod_{j=0}^{k-1} (x-j)
\end{equation}
is the falling factorial (some other notations have been used
for these polynomials, here we follow~\cite{Knuth} or \cite[\S\,2.6, p.~47]{GKP}). 
The discrete counterpart to the Bernoulli polynomials are the now so-called 
Bernoulli polynomials of the second kind (see~\cite{Car}),
that we will denote by $b_k(x)$, 
that were independently 
introduced by Jordan \cite{Jor} and Rey Pastor \cite{RP-BSMA} in 1929.
These polynomials, that have also been known with the name of Rey Pastor polynomials
(see~\cite{Bab}), are now defined by a generating function as in \eqref{eq:dAs} 
by means of 
\begin{equation}
\label{eq:fungenAn}
  \frac{t}{\log(1+t)}(1+t)^x = \sum_{k=0}^\infty b_k(x) \, \frac{t^k}{k!}.
\end{equation}
For $x=0$, the numbers $b_k(0)$ (or $b_k(0) \cdot k!$) appear, for instance, 
in Gregory's method for numerical integration (see~\cite{Ph}), and are also
called Bernoulli numbers of the second kind 
(see \cite[\S\,24.16]{Dil-NIST} or \cite{WuPan}), 
logarithmic numbers, Gregory coefficients, and 
Cauchy numbers of the first kind (see, for instance, \cite{MSV, Zh}).

There are many kind of generalizations of Bernoulli polynomials
that are defined by means of parameters in the function $A(t) = t/(e^t-1)$
(see, e.g., the classical papers
\cite{Ap, Nor} or the recent \cite{AdLe, NRV-Brno, Tem}). 
A complete different kind of generalization is obtained by replacing
the operator $\frac{d}{dx}$ in~\eqref{eq:A} instead of changing
the function~$A(t)$.

This was done in~\cite{CDPV}, where the Dunkl operator on the real line
(for the group $\mathbb{Z}_2$), 
namely
\begin{equation}
\label{eq:La}
  \La f(x) = \frac{d}{dx}f(x) + \frac{2\alpha+1}{2}
  \left(\frac{f(x)-f(-x)}{x}\right),
\end{equation}
where $\alpha > -1$ is a fixed parameter (see \cite{Du, Ros}),
was used instead of the ordinary derivative $\frac{d}{dx}$. 
In that setting, an Appell-Dunkl sequence $\{P_{k,\alpha}\}_{k=0}^{\infty}$ is
a sequence of polynomials that satisfies
\begin{equation}
\label{eq:LaPk}
  \La P_{k,\alpha}(x) = \big(k+(\alpha+1/2)(1-(-1)^k)\big) P_{k-1,\alpha}(x)
\end{equation}
(instead of $\La P_{k,\alpha} = k P_{k-1,\alpha}$, the previous definition uses 
another multiplicative constant in the place of $k$ for convenience with the notation). 
Of course, in the case $\alpha=-1/2$,
the operator $\La$ is the ordinary derivative and Appell-Dunkl sequences become
the classical Appell sequences. To give Appell-Dunkl sequences by means of 
a generating function, some extra notation is required.
Along all the paper, to clearly distinguish Dunkl cases from classical cases,
we will always include $\alpha$ as subindex in the notation of the polynomials
that come from the Dunkl case.

For $\alpha>-1$, we consider the entire function
\[
  \I_\alpha(z)
  = 2^\alpha \Gamma(\alpha+1)\, \frac{J_\alpha(iz)}{(iz)^\alpha}
  = \Gamma(\alpha+1) \sum_{n=0}^\infty \frac{(z/2)^{2n}}{n!\,\Gamma(n+\alpha+1)}
  = {}_0F_1(\alpha+1,z^2/4), 
\]
where $J_\alpha$ is the Bessel function of order~$\alpha$
(the function $\I_\alpha$ is a small variation of the so-called modified Bessel function of the first kind and order~$\alpha$, usually denoted by~$I_\alpha$, see~\cite{Wat} or~\cite{OlMa-NIST}),
and 
\[
  \Ea(z) = \I_\alpha(z) + \frac{z}{2(\alpha+1)} \, \I_{\alpha+1}(z),
  \qquad z \in \CC.
\]
From the very definition,
\[
  \Ea(z) = \sum_{n=0}^\infty \frac{z^n}{\cc_{n,\alpha}}
\]
with
\begin{equation}
\label{eq:ccna}
\cc_{n,\alpha} =
\begin{cases}
  2^{2k}k!\,(\alpha+1)_k, & \text{if $n=2k$},\\
  2^{2k+1}k!\,(\alpha+1)_{k+1}, & \text{if $n=2k+1$},
\end{cases}
\end{equation}
and where $(a)_n$ denotes the Pochhammer symbol
\[
  (a)_n = a(a+1)(a+2) \cdots (a+n-1) = \frac{\Gamma(a+n)}{\Gamma(a)}
\]
(with $n$ a non-negative integer). Notice that $\cc_{n,-1/2} = n!$ and $E_{-1/2}(z) = e^{z}$.
An important property is that, for any $\lambda \in \CC$, we have
\begin{equation}
\label{eq:LaEalx}
  \La \Ea(\lambda x) = \lambda E_\alpha(\lambda x),
\end{equation}
which is a generalization of $\frac{d}{dx} e^{\lambda x} = \lambda e^{\lambda x}$.

The function $\Ea(z)$ is known as the Dunkl kernel because, in a similar way to the Fourier transform (which takes place for $\alpha = -1/2$), 
we can define the Dunkl transform on the real line
\begin{equation}
\label{eq:trDu}
  \Fa f(y)
  = \int_{\RR} \Ea(-ixy)f(x)\,d\mua(x),
  \qquad y \in \RR,
\end{equation}
where $d\mua$ denotes the measure
\[
  d\mua(x) = \tfrac{1}{2^{\alpha+1}\Gamma(\alpha+1)}\, |x|^{2\alpha+1}\,dx
\]
(in particular, $d\mu_{-1/2}(x) = (2\pi)^{-1/2}\,dx$).
This operator has been widely studied in the mathematical literature
(see, for instance, \cite{Jeu, Ros, Rosler-LN, Tr, AdJ, NS, CPRV, DW, CRV, MO}); 
however, the goal of this paper is not the Dunkl transform, but some families of
polynomials that are defined with the aid of the Dunkl kernel~$\Ea$.
%The Dunkl transform~\eqref{eq:trDu} can be extended to an isometric isomorphism 
%on $L^2(\RR,\mua)$ and fulfils $\Fa^{-1}f(y) = \Fa f(-y)$.

As far as we are concerned, the entire function $\Ea$ 
is invariant under the Dunkl operator \eqref{eq:La}
%\begin{equation}
%\label{eq:La}
%  \La f(x) = \frac{d}{dx}f(x) + \frac{2\alpha+1}{2}
%  \left(\frac{f(x)-f(-x)}{x}\right)
%\end{equation}
in the same way that the exponential function is invariant under the ordinary derivative.
Then, it is easy to check that, if we have an analytic function $A(t)$
defined in a neighborhood of $0$ with $A(0) \ne 0$, and we take
\begin{equation}
\label{eq:APs}
  A(t)\Ea(xt) = \sum_{k=0}^\infty P_{k,\alpha}(x) \frac{t^k}{\cc_{k,\alpha}},
\end{equation}
then $P_{k,\alpha}(x)$ is a polynomial of degree $k$ which satisfies~\eqref{eq:LaPk}.
Thus, \eqref{eq:APs} is the characterization of Appell-Dunkl sequences
by means of a generating function.

The Appell-Dunkl polynomials had already appeared in the literature as 
generalizations of the Hermite polynomials (see for example \cite{AlS, Ros}).
The extension of Bernoulli polynomials to the Dunkl context was carried out in~\cite{CDPV};
there, Bernoulli-Dunkl polynomials $\{\B_{k,\alpha}\}_{k=0}^{\infty}$
are defined by a generating function via
\begin{equation}
\label{eq:defbd}
  \frac{\Ea(xt)}{\I_{\alpha+1}(t)}
  = \sum_{k=0}^\infty \B_{k,\alpha}(x) \frac{t^k}{\cc_{k,\alpha}},
\end{equation}
and it is shown that it represents a ``genuine extension'' because it is useful
to extend many standard properties of the classical Bernoulli polynomials
to the Dunkl context. In particular, how to use them 
to pose the series $\sum_{n=1}^{\infty} 1/n^{2k}$ in the Dunkl context
and to sum them. Some other Appell-Dunkl polynomials were studied
in \cite{DPV} and~\cite{CMV}.

The goal of this paper is to provide a suitable definition of discrete Appell-Dunkl
sequences, that is, to give an extension of~\eqref{eq:Adiscrete} and \eqref{eq:dAs} in the Dunkl context,
and to give some representative examples and properties. 
In particular, we show how to define Bernoulli-Dunkl polynomials of the second kind.

We need to notice an important obstacle. The ordinary translation $f(x) \mapsto f(x+1)$ is
useless in the Dunkl context, so the definition \eqref{eq:Adiscrete} must be adapted
taking it into account. 
There is, however, an analogue, the Dunkl translation operator $\tau_y$ which acts on a function $f$ by
\begin{equation}
\label{eq:tauyfx}
  \tau_y f(x) = \sum_{n=0}^\infty \frac{y^n}{\cc_{n,\alpha}}\La^{n}f(x),
  \qquad \alpha>-1,
\end{equation}
where $\La^0$ is the identity operator and $\La^{n+1} = \La(\La^n)$. 
When $\alpha=-1/2$, $\La^{n}$ is the $n$th derivative, 
so \eqref{eq:tauyfx} is the Taylor expansion of $f(x+y)$ around~$x$.
For the Dunkl transform, the translation $\tau_y$ plays the same role as the classical translation for the Fourier transform (that is, $\tau_y f(x) = f(x+y)$ for the case $\alpha=-1/2$). 
Two useful properties of the translation operator are the commutativity $\tau_a\tau_b = \tau_b\tau_a$ and the fact $\tau_a f(b) = \tau_b f(a)$.
Some other properties including an integral expression, can be found in \cite{Ros}, \cite{Rosler-TAMS}, and~\cite{ThX}.
In particular, let us note the identity \cite[formula (4.2.2)]{Ros}
\begin{equation}
\label{eq:Etxy}
  \tau_y(E_\alpha(t\cdot))(x) = E_\alpha(tx) E_\alpha (ty),
\end{equation}
that is a Dunkl alternative to the formula $e^{t(x+y)} = e^{tx}e^{ty}$ (case $\alpha=-1/2$).

Many properties of the Appell sequences of polynomials can be adapted to the Appell-Dunkl sequences
using the Dunkl translation operator.
For instance, in~\cite{CDPV, CMV} we can see that the Appell-Dunkl polynomials satisfy
\begin{equation}
\label{eq:tau-P}
  \tau_y(P_{k,\alpha}(\cdot))(x)
  = \sum_{j=0}^k \binoma{k}{j} P_{j,\alpha}(x) y^{k-j},
\end{equation}
with $\binoma{k}{j} = \cc_{k,\alpha}/(\cc_{j,\alpha}\cc_{k-j,\alpha})$;
in the classical case $\alpha=-1/2$, \eqref{eq:tau-P} becomes the well-known binomial formula
\begin{equation}
\label{eq:binom}
  P_k(x+y) = \sum_{j=0}^k \binom{k}{j} P_j(x)y^{k-j}
\end{equation}
for Appell polynomials. 
Some other properties concerning the use of Dunkl translations
for Appell-Dunkl polynomials can be found in~\cite{CMV}.


The structure of this paper is as follows. 
In Section~\ref{sec:DAD}, we see how the definition~\eqref{eq:Adiscrete} can be successfully adapted to the Dunkl context by using the Dunkl translation~\eqref{eq:tauyfx}, and what generating expansion plays the role of~\eqref{eq:dAs}.
In Section~\ref{sec:fallfact} we extend the classical falling factorial polynomials~\eqref{eq:ffact} to the Dunkl context, and we found their expansion by means of a discrete Appell-Dunkl generating function;
we will also see that these polynomials can be used to have a binomial formula that is analogous to \eqref{eq:binom} and \eqref{eq:tau-P} in the discrete Appell-Dunkl setting.
In Section~\ref{sec:another} we study some families of discrete Appell-Dunkl sequences than can be expressed in terms of the generalized Bernoulli-Dunkl polynomials that were defined in~\cite{CMV}. 
In Section~\ref{sec:connection} we show some relations between the families of polynomials defined in the Sections~\ref{sec:fallfact} and~\ref{sec:another}, and we prove some interesting results, in particular the Newton expansion in the Dunkl context.
In Section~\ref{sec:BDsecond} we define the Bernoulli-Dunkl polynomials of the second kind, as well as their corresponding generalization of order~$r$, and prove some properties.

%----------
\section{Discrete Appell-Dunkl polynomials}
\label{sec:DAD}
%----------

Let us start noticing that the Dunkl operator has a ``symmetric flavor'' due to the summand $(f(x)-f(-x))/x$ in \eqref{eq:La}; then, Appell-Dunkl polynomials must be defined in a symmetric way around $x=0$. However, the ``basic'' interval for the classical Bernoulli polynomials $\{B_k(x)\}_{k=0}^{\infty}$ is the interval $[0,1]$. In this way, Bernoulli-Dunkl polynomials are defined with ``basic'' interval $[-1,1]$, and a slight change of variable is needed to recover the classical Bernoulli polynomials in the case $\alpha=-1/2$. Particularly, we have
\begin{equation}
\label{eq:BtoB}
  \frac{\B_{k,-1/2}(2x-1)}{2^k} = B_k(x),
\end{equation}
see \cite{CDPV} or \cite{CMV} for details.

The same adjustment happens in the discrete case. The definition~\eqref{eq:Adiscrete} does not have a symmetric flavor, so it is not suitable for the Dunkl context. 
Then, instead of using the discrete operator $\Delta f(x) = f(x+1) - f(x)$ (forward differences), it is convenient to consider the classical case with the central discrete operator $\Delta f(x) = (f(x+1) - f(x-1))/2$ (central differences),
and the central discrete Appell sequences defined by means of the relation
\begin{equation}
\label{eq:Adiscen}
  \frac{p_k(x+1) - p_k(x-1)}{2} = k p_{k-1}(x), \quad k \ge 1.
\end{equation}
Although \eqref{eq:Adiscrete} is more usual in the classical case, discrete Appell sequences defined as in \eqref{eq:Adiscen} have also been studied; see for instance~\cite[\S\,6]{Tem}. 

In the same way that we have
\[
  f(x+1,t) - f(x,t) = t f(x,t)
\]
for $f(x,t) = (1+t)^x$, and this is the reason to the relation between \eqref{eq:Adiscrete} and~\eqref{eq:dAs}, we have that
\begin{equation}
\label{eq:eigendis}
  \frac{f(x+1,t) - f(x-1,t)}{2} = t f(x,t)
\end{equation}
for 
\begin{equation}
\label{eq:soleigendis}
  f(x,t) = (t+\sqrt{1+t^2})^x
  = \exp\big(x \log(t+\sqrt{1+t^2})\big).
\end{equation}
Then, the central discrete Appell sequences can also be defined by means of a generating function
\begin{equation}
\label{eq:dAsc}
  A(t)(t+\sqrt{1+t^2})^x = \sum_{k=0}^{\infty} p_k(x) \frac{t^k}{k!},
\end{equation}
where $A(t)$ is a function analytic at $t=0$ with $A(0) \ne 0$. 
(In particular, and among other examples, \cite[\S\,6]{Tem} defines the central Bernoulli polynomials of the second kind, we will see it in the forthcoming sections.)

In what follows, we will sometimes denote the Dunkl operator $\La$ and the Dunkl translation $\tau_a$ by $\Lax$ and $\tau_{a,x}$, respectively, to emphasize that the involved variable is~$x$. In fact, we will always use this notation when we apply these operators to a function with two variables. 
% $x$ and $t$, 
In addition, we will often use $\cc_k$ instead of $\cc_{k,\alpha}$.
% (defined in \eqref{eq:ccna})

To extend \eqref{eq:Adiscen} to the Dunkl context, that is, by using the Dunkl translation, it is enough to take sequences of polynomials that, for some constants $\theta_k$, satisfy the relation
\[
  (\alpha+1)(\tau_{1}-\tau_{-1}) p_{k,\alpha}(x) = \theta_k p_{k-1,\alpha}(x)
\]
(that, in the case $\alpha=-1/2$, recovers~\eqref{eq:Adiscen}).
But, which is the generating function? How can one find examples?
The key point is to give a suitable extension of \eqref{eq:eigendis} and~\eqref{eq:soleigendis} to the new scheme. We will often denote $\Delta_\alpha = (\alpha+1)(\tau_{1}-\tau_{-1})$, and we will use $\Delta_{\alpha,x}$ (as in the case of $\tau_{\pm1,x}$) if it is convenient to emphasize that the involved variable is~$x$.

In the Dunkl setting, the function $\Ea$ plays the role of the exponential function.
Then, let us take the equation
\begin{equation}
\label{eq:tau1-1f}
  (\alpha+1)(\tau_{1,x}-\tau_{-1,x}) f(x,t) = t f(x,t)
\end{equation}
and let us look for a solution of the form $f(x,t) = \Ea(x h(t))$; we want to find $h(t)$.

Recall that $\Ea(x h(t)) = \sum_{k=0}^{\infty} x^k h(t)^k/\cc_k$ and
$\Lax \Ea(x h(t)) = h(t) \Ea(x h(t))$. Then,
\begin{equation}
\label{eq:E(h(t))}
\begin{aligned} % No he escrito el (\alpha+1)
  (\tau_{1,x}-\tau_{-1,x}) \Ea(x h(t)) 
  &= \tau_{1,x} \Ea(x h(t)) - \tau_{-1,x} \Ea(x h(t)) \\
  &= \sum_{k=0}^{\infty} \frac{\Lax^k \Ea(x h(t))}{\cc_k}
  - \sum_{k=0}^{\infty} \frac{(-1)^k \Lax^k \Ea(x h(t))}{\cc_k} \\
  &= \sum_{k=0}^{\infty} \frac{h(t)^k \Ea(x h(t))}{\cc_k}
  - \sum_{k=0}^{\infty} \frac{(-1)^k h(t)^k \Ea(x h(t))}{\cc_k} \\
  &= 2 \sum_{k=0}^{\infty} \frac{h(t)^{2k+1} \Ea(x h(t))}{\cc_{2k+1}}.
\end{aligned}
\end{equation}
Therefore,
\[
  (\alpha+1)(\tau_{1,x}-\tau_{-1,x}) \Ea(x h(t)) = t \Ea(x h(t))
\]
becomes
\begin{equation}
\label{eq:h(t)}
  \sum_{k=0}^{\infty} \frac{h(t)^{2k+1}}{\cc_{2k+1}} = \frac{t}{2(\alpha+1)}.
\end{equation}
Here we have the odd terms of $\Ea(z) = \sum_{k=0}^{\infty} z^k/\cc_k$, and the odd part or $\Ea(z)$ is $\frac{z}{2(\alpha+1)}\I_{\alpha+1}(z)$.
Hence, \eqref{eq:h(t)} can be written as
$\frac{h(t)}{2(\alpha+1)}\I_{\alpha+1}(h(t)) = \frac{t}{2(\alpha+1)}$, or 
\begin{equation}
\label{eq:h(t)I}
h(t)\I_{\alpha+1}(h(t)) = t.
\end{equation}
Now, let us take 
\[
  G_{\alpha}(z) = z \I_{\alpha+1}(z) = z \,{}_0F_1(\alpha+2,z^2/4)
  = \Gamma(\alpha+1) \sum_{n=0}^\infty \frac{z^{2n+1}}{2^{2n} n!\,\Gamma(n+\alpha+1)}.
\]
This function is odd, non-negative for $z>0$, and increasing (for $z>0$, the derivative term by term of the series is positive),
so there exists the inverse function
\[
  h(t) = G_{\alpha}^{-1}(t), \quad t \in \RR,
\]
and this function satisfies~\eqref{eq:h(t)I}. 
This implies that the function 
\[
  f(x,t) = \Ea(x h(t)) = \Ea(x G_{\alpha}^{-1}(t))
\]
is the solution of~\eqref{eq:tau1-1f}.

Then, we would like to define a discrete Appell-Dunkl sequence as $\{p_{k,\alpha}(x)\}_{k=0}^{\infty}$ whose generating function is
\begin{equation}
\label{eq:defADd}
  A(t) \Ea(x G_{\alpha}^{-1}(t)) = \sum_{k=0}^{\infty} p_{k,\alpha}(x) \frac{t^k}{\cc_{k,\alpha}},
\end{equation}
where $A(t)$ is an analytic function in a neighborhood of $0$ satisfying $A(0) \ne 0$.

Let us now analyze the case $\alpha=-1/2$ for a moment.
It is easy to check that $\I_{1/2}(z) = \sinh(z)/z$, so 
% $G_{-1/2}(z) = \sinh(z)$ and
\[
  G_{-1/2}^{-1}(t) = \arcsinh(t) = \log(t+\sqrt{1+t^2}).
\]
Then, 
% using tat $E_{-1/2}=\exp$,
\[
  f(x,t) = E_{-1/2}(x G_{-1/2}^{-1}(t)) 
  = \exp(x \log(t+\sqrt{1+t^2})) = \big(t+\sqrt{1+t^2})^x,
\]
and, as expected, \eqref{eq:defADd} for $\alpha=-1/2$ becomes the classical central discrete Appell sequences that appear in~\eqref{eq:dAsc}.

For general $\alpha>-1$, the functions $\I_{\alpha}(z)$ are related to Bessel functions, and thus $G_{\alpha}^{-1}$ is the inverse of a function expressed in terms of them. As long as we know, this kind of inverse functions, that would be a generalization of the $\arcsinh$ function, has not been widely studied, and does not have a name. An example of a paper dealing with similar functions is~\cite{Gold}.
Later in this paper, we will give the analytic expansion of the function $G_{\alpha}^{-1}$ in terms of the generalized Bernoulli-Dunkl polynomials, see Theorem~\ref{teo:inver-Ga}.

Finally, let us state the main result of this section:

\begin{theorem}
\label{teo:defADdis}
Let $A(t)$ be an analytic function in a neighborhood of $0$ with $A(0) \ne 0$, and let $\{p_{k,\alpha}(x)\}_{k=0}^{\infty}$ be the sequence obtained by means of the generating function expansion
\begin{equation}
\label{eq:defADdis} 
  A(t) \Ea(x G_{\alpha}^{-1}(t)) = \sum_{k=0}^{\infty} p_{k,\alpha}(x) \frac{t^k}{\cc_{k,\alpha}},
  \qquad \alpha > -1.
\end{equation}
Then, $p_{k,\alpha}(x)$ is a polynomial of degree $k$, and it satisfies
\begin{equation}
\label{eq:tau1-1pk}
  (\alpha+1)(\tau_{1,x}-\tau_{-1,x}) p_{k,\alpha}(x) = \theta_{k,\alpha} p_{k-1,\alpha}(x),
  \qquad k=1,2,\ldots,
\end{equation}
with $\theta_{k,\alpha} = (k+(\alpha+1/2)(1-(-1)^k)) = \cc_{k,\alpha}/\cc_{k-1,\alpha}$.
\end{theorem}

\begin{proof}
Both functions $A(t)$ and $\Ea(x G_{\alpha}^{-1}(t))$ are analytic in a neighborhood of $t=0$, so the expansion \eqref{eq:defADdis} exists. 
Furthermore, $p_0(x) = A(0) \Ea(x G_{\alpha}^{-1}(0)) = A(0) \Ea(0) = A(0) \ne 0$, a constant.

For brevity, let us denote $\Delta_{\alpha,x} = (\alpha+1)(\tau_{1,x}-\tau_{-1,x})$;
and we will use $\theta_k$ and $\cc_k$ in the place of $\theta_{k,\alpha}$ and $\cc_{k,\alpha}$. 
By applying $\Delta_{\alpha,x}$ to both sides of \eqref{eq:defADdis}, we have
\begin{equation}
\label{eq:T1x}
  A(t) \Delta_{\alpha,x} \Ea(xG_{\alpha}^{-1}(t)) 
  = \sum_{k=0}^{\infty} \Delta_{\alpha,x} p_{k,\alpha}(x) \frac{t^k}{\cc_k}.
\end{equation}
By repeating the arguments in \eqref{eq:E(h(t))} and \eqref{eq:h(t)} with $G_{\alpha}(t) = t\I_{\alpha+1}(t)$, we get 
\begin{align*}
  \Delta_{\alpha,x} \Ea(xh(t)) 
  &= (\alpha+1)(\tau_{1,x}-\tau_{-1,x}) \Ea(x h(t)) 
  = 2(\alpha+1) \Ea(x h(t))  \sum_{k=0}^{\infty} \frac{h(t)^{2k+1}}{\cc_{2k+1}} \\
  &= \Ea(x h(t)) h(t) \I_{\alpha+1}(h(t))
  = \Ea(x h(t)) G_{\alpha}(h(t)).
\end{align*}
Then, taking $h(t) = G_{\alpha}^{-1}(t)$, the left hand side of \eqref{eq:T1x} is
\[
  \Delta_{\alpha,x} \Ea(x G_{\alpha}^{-1}(t)) = t \Ea(x G_{\alpha}^{-1}(t)).
\]
For the right hand side of \eqref{eq:T1x}, let us observe that $\tau_{\pm1,x}p_{0,\alpha}(x) = p_{0,\alpha}(x)$ (the translation of a constant is itself), so the term corresponding to $k=0$ vanishes. We can then simplify $t$ in both sides of \eqref{eq:T1x} so, taking  $l=k-1$, \eqref{eq:T1x} becomes
\[
  A(t) \Ea(xG_{\alpha}^{-1}(t)) 
  = \sum_{l=0}^{\infty} \Delta_{\alpha,x}p_{l+1,\alpha}(x) \frac{t^l}{\cc_{l+1}}.
\]
But $A(t) \Ea(xG_{\alpha}^{-1}(t)) = \sum_{l=0}^{\infty} p_{l,\alpha}(x) t^l/\cc_{l}$, and, since the power expansion of an analytic function is unique, $\Delta_{\alpha,x}p_{l+1,\alpha}(x) = \theta_{l+1} p_{l,\alpha}(x)$ with $\theta_{l+1} = \cc_{l+1}/\cc_l$.

Finally, taking into account that $p_{0,\alpha}(x)$ is constant (that is, a polynomial of degree~$0$), we easily get that $p_{k,\alpha}(x)$ is a polynomial of degree~$k$ by induction.
\end{proof}

%----------
%\section{The (falling) factorial polynomials in the Dunkl context}
\section{The factorial polynomials in the Dunkl context}
\label{sec:fallfact}
%----------

As explained in the introduction, the classical discrete case 
$A(t)(1+t)^x$ for $A(t)=1$ with forward differences
$\Delta f(x) = f(x+1)-f(x)$ gives
%  \eqref{eq:dAs}
\[
(1+t)^x = \sum_{k=0}^{\infty} x^{\underline{k}} \, \frac{t^k}{k!}
\]
where $x^{\underline{k}}$ is the falling factorial \eqref{eq:ffact}.

The classical case considering central differences $\Delta f(x) = (f(x+1)-f(x-1))/2$
and $A(t) = 1$ is
\[
(t+\sqrt{1+t^2})^x = \sum_{k=0}^{\infty} f_k(x) \, \frac{t^k}{k!}.
\]
In this case, the ``central factorial'' polynomials $f_k(x)$ have the following pattern:
\begin{equation}
\label{eq:clascentral}
\begin{aligned}
f_0(x) &= 1, \qquad
f_1(x) = x, \qquad
f_2(x) = x^2, \\
f_3(x) &= (-1+x) x (1+x), \qquad
f_4(x) = (-2+x) x^2 (2+x), \\
f_5(x) &= (-3+x) (-1+x) x (1+x) (3+x), \\
f_6(x) &= (-4+x) (-2+x) x^2 (2+x) (4+x), \\
f_7(x) &= (-5+x) (-3+x) (-1+x) x (1+x) (3+x) (5+x).
\end{aligned}
\end{equation}

For the Dunkl case, let us take \eqref{eq:defADd} with $A(t)=1$.
Then, we can say that the Dunkl factorial (or Dunkl ``central'' factorial)
% (we will drop the word ``falling'' because it is a case of ``central factorial'') 
are the polynomials  
$\{f_{k,\alpha}(x)\}_{k=0}^{\infty}$ whose generating function is
\begin{equation}
\label{eq:fallfka}
  \Ea(x G_{\alpha}^{-1}(t)) = \sum_{k=0}^{\infty} f_{k,\alpha}(x) \frac{t^k}{\cc_{k,\alpha}}.
\end{equation}
It is not difficult to check that the first Dunkl factorial polynomials are
\begin{align*}
f_{0,\alpha}(x) &= 1, \qquad
f_{1,\alpha}(x) = x, \qquad
f_{2,\alpha}(x) = x^2, \\
f_{3,\alpha}(x) &= x^3-x, \qquad
f_{4,\alpha}(x) = x^4 - 4x^2, \\
f_{5,\alpha}(x) &= x^5
  - \frac{6(\alpha+3) x^3}{\alpha+2}
  + \frac{(5\alpha+16) x}{\alpha+2}, \\
f_{6,\alpha}(x) &= x^6
  - \frac{12(\alpha+3)x^4}{\alpha+2}
  + \frac{6(6\alpha+19) x^2}{\alpha+2} \\
f_{7,\alpha}(x) &= x^7
  - \frac{15(\alpha+4) x^5}{\alpha+2}
  + \frac{9(\alpha+4)(7\alpha+22) x^3}{(\alpha+2)^2}
  - \frac{(7\alpha+26)^2 x}{(\alpha+2)^2}.
\end{align*}


It is perhaps surprising that the polynomials $f_{k,\alpha}$ do not have any 
recognizable pattern as in \eqref{eq:clascentral}.
But the same happens when the falling factorials are defined in other contexts;
indeed, this is what happens in \cite[\S\,18]{Cho},
where they are called factor polynomials
(a more detailed explanation of the similarities and the differences
between the context in~\cite{Cho} and our context can be found 
in \cite[Remark~1]{CMV}).

In the discrete Appell context, there is a binomial formula that is the discrete alternative to the binomial formula~\eqref{eq:binom}: 
if $\{p_k(x)\}_{k=0}^{\infty}$ is a discrete Appell
sequence (that is, it is defined as in \eqref{eq:dAs}), they satisfy
\begin{equation}
\label{eq:binomdis}
  p_k(x+y) = \sum_{j=0}^k \binom{k}{j} p_j(x)y^{\underline{k-j}}.
\end{equation}
Let us see how to adapt this expression in the discrete Appell-Dunkl setting, that is, how to obtain a discrete version of~\eqref{eq:tau-P}.
Notice that the role of the factorial polynomials $y^{\underline{k}}$ in~\eqref{eq:binomdis} is now played by their Dunkl counterpart $f_{k,\alpha}(y)$.

\begin{theorem}
Let $\alpha>-1$ and $\{p_{k,\alpha}(x)\}_{k=0}^{\infty}$ be a discrete Appell-Dunkl sequence of polynomials 
defined as in \eqref{eq:defADdis}. Then, we have
\[
   \tau_y(p_{k,\alpha}(\cdot))(x)
  = \sum_{j=0}^k \binoma{k}{j} p_{j,\alpha}(x) f_{k-j,\alpha}(y),
\]
where $\{f_{k,\alpha}(y)\}_{k=0}^{\infty}$ are the Dunkl factorial polynomials defined in~\eqref{eq:fallfka}.
\end{theorem}

\begin{proof}
By using~\eqref{eq:Etxy}, we have
\[
  \tau_y(\Ea(\,\cdot\, G_\alpha^{-1}(t))(x) 
  = \Ea(xG_\alpha^{-1}(t)) \Ea (yG_\alpha^{-1}(t)).
\]
Then, if we apply $\tau_y$ to \eqref{eq:defADdis} we get
\begin{align*}
  &\sum_{k=0}^{\infty} \tau_y(p_{k,\alpha})(x) \frac{t^k}{\cc_{k}}
  = A(t) \tau_y(\Ea(\,\cdot\, G_\alpha^{-1}(t))(x)
  = A(t) \Ea(xG_\alpha^{-1}(t)) \Ea (yG_\alpha^{-1}(t)) \\
  &\qquad= \bigg(\sum_{k=0}^{\infty} p_{k,\alpha}(x) \frac{t^k}{\cc_{k}}\bigg)
  \bigg(\sum_{k=0}^{\infty} f_{k,\alpha}(y) \frac{t^k}{\cc_{k}}\bigg) 
  = \sum_{k=0}^{\infty} \bigg( \sum_{j=0}^k \frac{p_{j,\alpha}(x)}{\cc_{j}} 
    \frac{f_{k-j,\alpha}(y)}{\cc_{k-j}} \bigg) t^k \\
  &\qquad= \sum_{k=0}^{\infty} \bigg( \sum_{j=0}^k \binoma{k}{j} p_{j,\alpha}(x) 
    f_{k-j,\alpha}(y) \bigg) \frac{t^k}{\cc_{k}}.
\end{align*}
By equating coefficients of $t^k$, we get the result.
\end{proof}

Let us finally see that, in a similar way to what happens in the traditional settings, the role played by the monomials $x^k$ in the Appell-Dunkl case is assumed by the factorial polynomials $f_{k,\alpha}(x)$ in the discrete Appell-Dunkl case. In \cite{CDPV, CMV} we saw that, if $\{P_{k,\alpha}(x)\}_{k=0}^{\infty}$ is a Appell-Dunkl sequence defined as in \eqref{eq:APs} with $1/A(t) = \sum_{k=0}^\infty a_k t^k$, we have
\[
  x^k = \cc_{n,\alpha} \sum_{j=0}^k \frac{P_{j,\alpha}(x)}{\cc_{j,\alpha}}\, a_{k-j}.
\]
For discrete Appell-Dunkl sequences, the corresponding result is as follows:


\begin{theorem}
\label{teo:fnfrompj}
Let $\alpha>-1$ and $\{p_{k,\alpha}(x)\}_{k=0}^{\infty}$ be a discrete Appell-Dunkl sequence of polynomials 
defined as in \eqref{eq:defADdis}, with $1/A(t) = \sum_{k=0}^\infty a_k t^k$. Then, we have
\begin{equation*}
%\label{eq:fnfrompj}
  f_{k,\alpha}(x) = \cc_{k,\alpha} \sum_{j=0}^k \frac{p_{j,\alpha}(x)}{\cc_{j,\alpha}}\, a_{k-j}.
\end{equation*}
\end{theorem}

\begin{proof}
Let us write \eqref{eq:defADdis} as
\[
\Ea(xG_{\alpha}^{-1}(t)) 
= \bigg( \sum_{k=0}^{\infty} p_{k,\alpha}(x) \frac{t^k}{\cc_k} \bigg) \frac{1}{A(t)}
\]
and express all the functions as their series,
\[
\sum_{k=0}^{\infty} f_{k,\alpha}(x) \frac{t^k}{\cc_k}
= \bigg( \sum_{k=0}^{\infty} p_{k,\alpha}(x) \frac{t^k}{\cc_k} \bigg)
\bigg( \sum_{k=0}^{\infty} a_k t^k \bigg)
= \sum_{k=0}^{\infty} \bigg( \sum_{j=0}^k \frac{p_{j,\alpha}(x)}{\cc_{j}}\, a_{k-j} \bigg) t^k.
\]
The result follows by identifying coefficients.
\end{proof}

%----------
\section{Discrete Appell-Dunkl families defined in terms 
of generalized Bernoulli-Dunkl polynomials}
\label{sec:another}
%----------

For an arbitrary number $r$ 
(we are mainly interested in the case of a non-negative integer, but it is not a restriction), 
the generalized Bernoulli polynomials of order $r$ (also known as N{\o}rlund polynomials) were defined by the first time in~\cite{Nor}, and they satisfy
\[
  \bigg(\frac{t}{e^t-1}\bigg)^r e^{xt} = \sum_{k=0}^{\infty} B_k^{(r)}(x) \frac{t^k}{k!}.
\]
We point out that the case $r=1$ corresponds to the ordinary Bernoulli polynomials, and the trivial case $r=0$ corresponds to the polynomials $\{x^k\}_{k=0}^{\infty}$.
It is well known that the polynomials $B_k^{(r)}(x)$ satisfy 
\begin{equation}
\label{eq:Bkrder}
  B_{k}^{(r+1)}(x+1) - B_{k}^{(r+1)}(x)
  = \frac{d}{dx} B_{k}^{(r)}(x) = k B_{k-1}^{(r)}(x)
\end{equation}
(see, for instance, \cite{Car}, \cite[\S\S\,6.11 and 6.3]{MT} or~\cite{Nor}).
Then, if we set 
\begin{equation}
\label{eq:PkDesdeBr}
  p_k(x) = B_{k}^{(k+1)}(x), \quad k=0,1,2,\dots,
\end{equation}
we have $p_k(x+1) - p_k(x) = k p_{k-1}(x)$, so \eqref{eq:PkDesdeBr} is an example of a discrete Appell sequence. Actually, also $\{B_{k}^{(k+s)}(x)\}_{k=0}^{\infty}$ is a discrete Appell sequence, but the case $s=1$ that has been used in~\eqref{eq:PkDesdeBr} is the most usual because the polynomials $B_k^{(k+1)}(x)$ have a well recognizable pattern (see \cite[\S\,6.4, p.~130]{MT} or \cite{Nor}), namely $B_0^{(1)}(x)=1$ and
\begin{equation}
\label{eq:Bkk+1}
  B_k^{(k+1)}(x) = (x-1)(x-2) \cdots (x-k) = (x-1)^{\underline{k}}
\end{equation}
(otherwise, even in the case $s=0$, $B_k^{(k+s)}(x)$ does not have an easy factorizable expression).
The generating function for the polynomials $p_k(x) = (x-1)^{\underline{k}}$ in~\eqref{eq:PkDesdeBr} is
\[
  \frac{1}{1+t} (1+t)^{x} = (1+t)^{x-1} = \sum_{k=0}^{\infty} (x-1)^{\underline{k}} \,\frac{t^k}{k!}.
\]

This can also be done in the classical central case (the details are left to the reader). We can define the generalized central Bernoulli polynomials and proceed as in~\eqref{eq:PkDesdeBr}. Then, we get the polynomials 
\begin{equation}
\label{eq:disgenc}
\begin{aligned}
\mathbf{p}_{0}(x) &= 1, \qquad
\mathbf{p}_{1}(x) = x, \qquad
\mathbf{p}_{2}(x) = (-1 + x) (1 + x), \\
\mathbf{p}_{3}(x) &= (-2 + x) x (2 + x), \qquad
\mathbf{p}_{4}(x) = (-3 + x) (-1 + x) (1 + x) (3 + x), \\
\mathbf{p}_{5}(x) &= (-4 + x) (-2 + x) x (2 + x) (4 + x), \\
\mathbf{p}_{6}(x) &= (-5 + x) (-3 + x) (-1 + x) (1 + x) (3 + x) (5 + x), \\
\mathbf{p}_{7}(x) &= (-6 + x) (-4 + x) (-2 + x) x (2 + x) (4 + x) (6 + x),
\end{aligned}
\end{equation}
that is a central discrete Appell sequence 
(i.e., they satisfy $\mathbf{p}_{k}(x+1)-\mathbf{p}_{k}(x-1) = 2k\mathbf{p}_{k-1}(x)$)
whose generating function is
\begin{equation}
\label{eq:disgenAc}
  \frac{1}{\sqrt{1+t^2}} \big(t + \sqrt{1+t^2}\big)^x 
  = \sum_{k=0}^{\infty} \mathbf{p}_k(x) \frac{t^k}{k!}.
\end{equation} 

A definition as in \eqref{eq:PkDesdeBr} can also be done in the Dunkl setting. In \cite{CMV}, the generalized Bernoulli-Dunkl polynomials $\{\B_{k,\alpha}^{(r)}(x)\}_{k=0}^{\infty}$ of order $r$ (where $r$ is an arbitrary real or complex parameter) are defined by means of the generating function
\begin{equation}
\label{eq:defbdr}
  \frac{\Ea(xt)}{\I_{\alpha+1}(t)^r}
  = \sum_{k=0}^\infty \B_{k,\alpha}^{(r)}(x) \frac{t^k}{\cc_{k,\alpha}}.
\end{equation}
In addition (see \cite[Theorem 8.2]{CMV}), they satisfy
\begin{equation}
\label{eq:tau1-1-Br}
  (\alpha+1)\big(\tau_1 \B_{k,\alpha}^{(r+1)}(x) - \tau_{-1} \B_{k,\alpha}^{(r+1)}(x)\big)
  = \La(\B_{k,\alpha}^{(r)})(x) = \theta_{k,\alpha} \B_{k-1,\alpha}^{(r)}(x).
\end{equation}
Then, if we take
\begin{equation}
\label{eq:PnDesdeBrD}
  \mathbf{p}_{k,\alpha}(x) = \B_{k,\alpha}^{(k+1)}(x), \quad k=0,1,2,\dots,
\end{equation}
we have $(\alpha+1)(\tau_{1,x}-\tau_{-1,x}) \mathbf{p}_{k,\alpha}(x) 
= \theta_{k,\alpha} \mathbf{p}_{k-1,\alpha}(x)$,
so $\{\B_{k,\alpha}^{(k+1)}(x)\}_{k=0}^{\infty}$ is a discrete Appell-Dunkl sequence.

Which is the generating function of this sequence? 
In the view of \eqref{eq:disgenAc}, one can guess that the generating function is \eqref{eq:defADdis} with $A(t) = (1+t^2)^{\alpha}$, but it is not true. 
In particular, this shows that a mere extension of the classical case to a more general with a parameter $\alpha$ in such a way that the value $\alpha=-1/2$ recovers the classical case is not always a good Dunkl extension. 

Actually, what happens is the following:

\begin{theorem}
\label{teo:Bkr}
Let $\alpha>-1$ and $\B_{k,\alpha}^{(r)}(x)$ denote the generalized Bernoulli-Dunkl polynomials defined by means of~\eqref{eq:defbdr}. Then $\{\B_{k,\alpha}^{(k+s)}(x)\}_{k=0}^{\infty}$ is a discrete Appell-Dunkl sequence whose generating function is
\begin{equation}
\label{eq:Bkrs}
  \bigg(\frac{t}{G_{\alpha}^{-1}(t)}\bigg)^{1-s} 
  \Big(\frac{d}{dt}(G_{\alpha}^{-1}(t))\Big) 
  \Ea(x G_{\alpha}^{-1}(t)) 
  = \sum_{k=0}^{\infty} \B_{k,\alpha}^{(k+s)}(x) \frac{t^k}{\cc_{k,\alpha}}, \qquad x \in \RR.
\end{equation}
In particular, for $s=1$ we have
\begin{equation}
\label{eq:Bkr}
  \Big(\frac{d}{dt}(G_{\alpha}^{-1}(t))\Big) \Ea(x G_{\alpha}^{-1}(t)) 
  = \sum_{k=0}^{\infty} \mathbf{p}_{k,\alpha}(x) \frac{t^k}{\cc_{k,\alpha}}, \qquad x \in \RR.
\end{equation}
\end{theorem}

\begin{proof}
We have from \eqref{eq:tau1-1-Br} that $\{\B_{k,\alpha}^{(k+s)}(x)\}_{k=0}^{\infty}$ is a discrete Appell-Dunkl sequence. Let us prove that \eqref{eq:Bkrs} is the generating function.

It is well known that the zeros of $J_{\alpha+1}(x)/x^{\alpha+1}$ are $\{s_{j,\alpha}\}_{j\in\ZZO}$, with $s_{j,\alpha} > 0$ and $s_{j+1,\alpha} > s_j$ for $j > 0$, and $s_{-j,\alpha}=-s_{j,\alpha}$ (see \cite[Chapter~15]{Wat} or \cite[\S\,10.21]{OlMa-NIST}). Moreover,
\[
  \I_{\alpha+1}(z)
  = 2^{\alpha+1} \Gamma(\alpha+2)\, \frac{J_{\alpha+1}(iz)}{(iz)^{\alpha+1}}
\]
so the function ${\Ea(xt)}/{\I_{\alpha+1}(t)^r}$ (as a function of the variable~$t$) is analytic in the disk $D(0,s_{1,\alpha})$ around $t=0$.
Then, by Cauchy's integral formula for the derivatives applied to~\eqref{eq:defbdr}, we have
\[
\frac{\B_{k,\alpha}^{(r)}(x)}{\cc_k}
= \frac{1}{2\pi i} \int_{C_1} \frac{\Ea(xt)}{\I_{\alpha+1}(t)^r} \frac{1}{t^{k+1}}\,dt,
\]
where $C_1$ is any circle centered at the origin with radius less than $s_{1,\alpha}$. 

The function $G_{\alpha}(t) = t\I_{\alpha+1}(t)$ is analytic in a neighborhood of $t=0$, and satisfies $G_{\alpha}(0)=0$ and $G_{\alpha}'(0)\ne 0$. Then, there exists $G_{\alpha}^{-1}(z)$ analytic in a neighborhood of $z=0$, and, for $C_1$ of radius small enough, $C_2 = G_{\alpha}(C_1)$ is a simple closed curve surrounding $z=0$. Thus, with the change $t = G_{\alpha}^{-1}(z)$ we have
\[
\frac{\B_{k,\alpha}^{(r)}(x)}{\cc_k}
= \frac{1}{2\pi i} \int_{C_2} \frac{\Ea(xG_{\alpha}^{-1}(z))}{z^r}
\frac{\frac{d}{dz}(G_{\alpha}^{-1}(z))}{G_{\alpha}^{-1}(z)^{k+1-r}}\,dz.
\]
Taking $r=k+s$, 
\[
\frac{\B_{k,\alpha}^{(k+s)}(x)}{\cc_k}
= \frac{1}{2\pi i} \int_{C_2} \frac{\Ea(xG_{\alpha}^{-1}(z))}{z^{s-1}(G_{\alpha}^{-1}(z))^{1-s}}
\frac{d}{dz}(G_{\alpha}^{-1}(z))\,\frac{dz}{z^{k+1}},
\]
and consequently 
\[
\bigg(\frac{t}{G_{\alpha}^{-1}(t)}\bigg)^{1-s} \Ea(xG_{\alpha}^{-1}(t)) 
\frac{d}{dt}G_{\alpha}^{-1}(t)
= \sum_{k=0}^{\infty}\frac{\B_{k,\alpha}^{(k+s)}(x)}{\cc_k} \,t^k.
\qedhere
\]
\end{proof}

Finally, let us note that the first polynomials $\mathbf{p}_{k,\alpha}(x) = \B_{k,\alpha}^{(k+1)}(x)$ are
\begin{align*}
\B_{0,\alpha}^{(1)}(x) &= 1, \qquad
\B_{1,\alpha}^{(2)}(x) = x, \qquad
\B_{2,\alpha}^{(3)}(x) = x^2 - \frac{3 (\alpha +1)}{\alpha +2}, \\
\B_{3,\alpha}^{(4)}(x) &= x^3 - 4 x, \qquad
\B_{4,\alpha}^{(5)}(x) = x^4 - 10 x^2 
  + \frac{5 (\alpha +1) (5 \alpha +16)}{(\alpha +2) (\alpha +3)}, \\
\B_{5,\alpha}^{(6)}(x) &= x^5 - \frac{12 (\alpha +3) x^3}{\alpha +2}
  + \frac{6 (6\alpha +19) x}{\alpha +2}, \\
\B_{6,\alpha}^{(7)}(x) &= x^6 - \frac{21 (\alpha +3) x^4}{\alpha +2}
  + \frac{21 (7 \alpha +22) x^2}{\alpha +2}
  - \frac{7 (\alpha +1) (7 \alpha +26)^2}{(\alpha +2)^2 (\alpha +4)}, \\
\B_{7,\alpha}^{(8)}(x) &= x^7 -\frac{24 (\alpha +4) x^5}{\alpha +2}
  + \frac{24 (\alpha +4) (8\alpha +25) x^3}{(\alpha +2)^2} 
  - \frac{32 (2\alpha+7)(8\alpha+31) x}{(\alpha +2)^2}.
\end{align*}

Of course, \eqref{eq:disgenc} corresponds to the case $\alpha=-1/2$  but, for general $\alpha>-1$, they cannot be factorized in a similar way.


%----------
%\section{Relation between the polynomials of Sections~\ref{sec:fallfact} 
%and~\ref{sec:another} and consequences}
\section{Relationships between families and applications}
\label{sec:connection}
%----------


The classical falling factorial polynomials $x^{\underline{k}}$ and the generalized Bernoulli polynomials $B_k^{(r)}$ are related by means of the formula
\[
    x^{\underline{k}} = \frac{x}{k} \frac{d}{dx} B_{k}^{(k)}(x),
    \quad k \ge 1.
\]
% En \cite[p.~130]{MT} y \cite[\S\,6.41, p.~133]{MT}, \cite{Car} hay f\'ormulas parecidas 
This clearly follows from $x^{\underline{k}} = xB_{k-1}^{(k)}$ (see~\eqref{eq:Bkk+1}) using~\eqref{eq:Bkrder}.

In the Dunkl context, we have the following:

\begin{theorem}
\label{teo:relation}
The Dunkl factorial polynomials $f_{k,\alpha}(x)$ defined in~\eqref{eq:fallfka}, and the generalized Bernoulli-Dunkl polynomials defined in~\eqref{eq:defbdr} are related by
\begin{equation}
\label{eq:relation}
    f_{k,\alpha}(x) = \frac{x}{k} \frac{d}{dx}\B_{k,\alpha}^{(k)}(x),
    \quad k = 1,2,\ldots.
\end{equation}
\end{theorem}

\begin{proof}
Let us check that the sequences of polynomials $\{kf_{k,\alpha}(x)\}_{k=0}^{\infty}$ and $\{x\frac{d}{dx}\B_{k,\alpha}^{(k)}(x)\}_{k=0}^{\infty}$ have the same generating function. 

For $\{kf_{k,\alpha}(x)\}_{k=0}^{\infty}$ we have, from \eqref{eq:fallfka}, that
\[
t \frac{d}{dt} \Big(\Ea(xG_{\alpha}^{-1}(t))\Big)
= t\sum_{k=0}^{\infty} \frac{f_{k,\alpha}(x)}{\cc_k}\,kt^{k-1} 
= \sum_{k=1}^{\infty} \frac{kf_{k,\alpha}(x)}{\cc_k}\,t^{k}.
\]
Moreover, let us notice that 
\[
\frac{d}{dt} \big(\Ea(xG_{\alpha}^{-1}(t))\big)
= \Ea'(xG_{\alpha}^{-1}(t)) x \frac{d}{dt} (G_{\alpha}^{-1}(t))
\]
where $\Ea'$ is the ordinary derivative of the Dunkl kernel $\Ea$. 
Then, we have the generating function
\begin{equation}
\label{eq:gf1}
tx \Ea'(xG_{\alpha}^{-1}(t)) \frac{d}{dt} (G_{\alpha}^{-1}(t))
= \sum_{k=1}^{\infty} \frac{kf_{k,\alpha}(x)}{\cc_k}\,t^{k}.
\end{equation}

On the other hand, let us use \eqref{eq:Bkrs} with $s=0$, that is,
\[
  \frac{t}{G_{\alpha}^{-1}(t)} 
  \Big(\frac{d}{dt}(G_{\alpha}^{-1}(t))\Big) 
  \Ea(x G_{\alpha}^{-1}(t)) 
  = \sum_{k=0}^{\infty} \B_{k,\alpha}^{(k)}(x) \frac{t^k}{\cc_{k}}.
\]
By differentiating this expression with respect to $x$, we have
\[
  t \Big(\frac{d}{dt}(G_{\alpha}^{-1}(t))\Big) 
  \Ea'(x G_{\alpha}^{-1}(t)) 
  = \sum_{k=0}^{\infty} \frac{d}{dx} \B_{k,\alpha}^{(k)}(x) \frac{t^k}{\cc_{k}},
\]
so
\begin{equation}
\label{eq:gf2}
  tx \Big(\frac{d}{dt}(G_{\alpha}^{-1}(t))\Big) 
  \Ea'(x G_{\alpha}^{-1}(t)) 
  = \sum_{k=0}^{\infty} x \frac{d}{dx} \B_{k,\alpha}^{(k)}(x) \frac{t^k}{\cc_{k}}.
\end{equation}
The generating functions in \eqref{eq:gf1} and \eqref{eq:gf2} are the same, so the result follows.
\end{proof}

In the place of \eqref{eq:relation}, another expansion can be found for $f_{k,\alpha}(x)$, also in terms of generalized Bernoulli polynomials.

We showed in \eqref{eq:tau-P} that every Appell-Dunkl sequence $\{P_{k,\alpha}(x)\}_{k=0}^{\infty}$
satisfies the binomial property 
\[
  P_{k,\alpha}(x) = \sum_{j=0}^k \binoma{k}{j} P_{j,\alpha}(0) x^{k-j}
\]
(recall that the Dunkl translation satisfies $\tau_y f(x) = \tau_x f(y)$, and that $\tau_0$ is the identity operator). Then, for the generalized Bernoulli-Dunkl polynomials we have
\begin{equation*}
%\label{eq:Bergen}
\B_{k,\alpha}^{(r)}(x) 
= \sum_{j=0}^k \binoma{k}{j} \B_{j,\alpha}^{(r)}(0) x^{k-j}.
\end{equation*}
If we differentiate this expression with respect to $x$, we obtain
\begin{equation}
\label{eq:derBka}
\frac{d}{dx} \B_{k,\alpha}^{(r)}(x) 
= \sum_{j=0}^{k-1} \binoma{k}{j} \B_{j,\alpha}^{(r)}(0)(k-j)x^{k-j-1}.
\end{equation}
Using together \eqref{eq:relation} and \eqref{eq:derBka} with $r=k$, we get
\begin{equation}
\label{eq:fkaBka}
f_{k,\alpha}(x) = \sum_{j=0}^{k-1} \binoma{k}{j} \B_{j,\alpha}^{(k)}(0) \Big(1-\frac{j}{k}\Big)x^{k-j}.
\end{equation}

This expression allows us to give the analytic expansion for the function~$G_\alpha^{-1}$:

\begin{theorem}
\label{teo:inver-Ga}
For $\alpha>-1$, the function $G^{-1}_\alpha$, inverse of $G_\alpha(z) = z \I_{\alpha+1}(z)$, can be given as
\[
G_{\alpha}^{-1}(t) 
= \sum_{k=1}^{\infty} \frac{\theta_{k}}{k} \frac{\B_{k-1,\alpha}^{(k)}(0)}{\cc_{k,\alpha}}t^k
= t \sum_{k=0}^{\infty} \frac{\B_{k,\alpha}^{(k+1)}(0)}{k+1} \frac{t^k}{\cc_{k,\alpha}}.
\]
\end{theorem}

\begin{proof}
Let us apply the Dunkl operator $\La$ to the function $\Ea(xG_{\alpha}^{-1}(t))$. 
Using \eqref{eq:LaEalx} and evaluating at $x=0$ gives
\[
\Lax(\Ea(xG_{\alpha}^{-1}(t)))\big|_{x=0} = G_{\alpha}^{-1}(t).
\]
Then, from \eqref{eq:fallfka} we have
\begin{equation}
\label{eq:G-1afka0}
G_{\alpha}^{-1}(t) = \sum_{k=0}^{\infty} \La f_{k,\alpha}(0) \frac{t^k}{\cc_k}.
\end{equation}

By applying $\La$ to \eqref{eq:fkaBka}, and using that $\La x^j = \theta_j x^{j-1}$, we get
\[
\La f_{k,\alpha}(x) 
= \sum_{j=0}^{k-1} \binoma{k}{j} \B_{j,\alpha}^{(k)}(0) \Big(1-\frac{j}{k}\Big)
\theta_{k-j}x^{k-j-1}.
\]
Evaluating at $x=0$, all the summands vanish except $j=k-1$, so
\[
\La f_{k,\alpha}(0) = \binoma{k}{k-1} \B_{k-1,\alpha}^{(k)}(0) \Big(1-\frac{k-1}{k}\Big)\theta_{1}
= \frac{\theta_k}{k} \B_{k-1,\alpha}^{(k)}(0).
\]
Then, \eqref{eq:G-1afka0} proves the theorem.
\end{proof}

Let us finish this section giving an extension of the Newton expansion to the Dunkl setting.
The Newton expansion consists of the terms of the Newton forward difference equation and,
in essence, it is the Newton interpolation formula, first published 
in his \textit{Principia Mathematica} in 1687. It is the discrete analog 
of the continuous Taylor expansion. 

To our concerns (see, for instance \cite[\S\,34, p.~191]{Nor}), the classical formula is the expansion of a polynomial $Q_n(x)$ of degree less than or equal to $n$ in terms of the falling factorials $x^{\underline{k}}$, namely
\begin{equation}
\label{eq:Qnx}
  Q_n(x) = \sum_{k=0}^n \frac{\Delta^kQ_n(0)}{k!}\, x^{\underline{k}}
\end{equation}
with $\Delta^{0} f(x) = f(x)$, $\Delta f(x) = f(x+1)-f(x)$ and $\Delta^{k+1} f(x) = \Delta (\Delta^k f(x))$. 

In the Dunkl case, we must use, instead of forward differences, 
the operator $(\alpha+1)(\tau_{1,x}-\tau_{-1,x})$;
and, instead of falling factorials $x^{\underline{k}}$, 
their Dunkl counterpart $f_{k,\alpha}(x)$. 
Then, the Newton-Dunkl expansion formula is the following:

\begin{theorem}
Let $\alpha>-1$ and $f_{k,\alpha}(x)$ be the Dunkl factorial polynomials defined in~\eqref{eq:fallfka}. 
Any polynomial $Q_n(x)$ of degree less than or equal to $n$ can be written as
\begin{equation}
\label{eq:QnxD}
  Q_n(x) = \sum_{k=0}^n \frac{\Delta_{\alpha}^kQ_n(0)}{\cc_{k,\alpha}}\, f_{k,\alpha}(x),
\end{equation} 
where $\Delta_{\alpha}^k$ is the operator $\Delta_{\alpha} = (\alpha+1)(\tau_{1}-\tau_{-1})$ applied $k$ times.
\end{theorem}

\begin{proof}
We know that $f_{0,\alpha}(x) = 1$ and 
$f_{k,\alpha}(x) = \frac{x}{k} \frac{d}{dx}\B_{k,\alpha}^{(k)}(x)$ 
(Theorem~\ref{teo:relation}). Consequently,
\begin{equation}
\label{eq:fk0}
f_{0,\alpha}(0) = 1 \quad\text{and}\quad f_{k,\alpha}(0) = 0 \text{ for }k \ge 1.
\end{equation}
The polynomials $\{f_{k,\alpha}(x)\}_{k=0}^n$ form a basis of the polynomials of degree $\le n$, so $Q_n(x)$ can be written as
\begin{equation}
\label{eq:Qn}
Q_n(x) = c_0 f_{0,\alpha}(x) + c_1 f_{1,\alpha}(x) + c_2 f_{2,\alpha}(x) 
+ \cdots + c_n f_{n,\alpha}(x)
\end{equation}
for some coefficients $c_k$. To compute the $c_k$, let us successively apply the operator $\Delta_{\alpha}$ to \eqref{eq:Qn}, and evaluate it at $x=0$.

In the first place, taking $x=0$ in \eqref{eq:Qn} gives $Q_n(0)=c_0$.
Secondly, applying $\Delta_{\alpha}$ to \eqref{eq:Qn} 
and using $\Delta_{\alpha}f_{k,\alpha} =  \theta_k f_{k-1,\alpha}$, gives
\[
\Delta_{\alpha} Q_n(x) = c_1 \theta_1 f_{0,\alpha}(x) + c_2 \theta_2 f_{1,\alpha}(x) 
+ \cdots + c_n \theta_n f_{n-1,\alpha}(x);
\]
and evaluating at $x=0$ gives $\Delta_{\alpha} Q_n(0) = c_1 \theta_1 = c_1 \cc_1$ (recall~\eqref{eq:fk0}), so
%\[
%c_1 = \frac{\Delta_{\alpha}Q_n(0)}{\cc_1}.
%\]
$c_1 = {\Delta_{\alpha}Q_n(0)}/{\cc_1}$.
Applying $\Delta_{\alpha}$ again we have
\[
\Delta_{\alpha}^2 Q_n(x) = c_2 \theta_2\theta_1 f_{0,\alpha}(x) + c_3 \theta_3\theta_2 f_{1,\alpha}(x) 
+ \cdots + c_n \theta_n\theta_{n-1} f_{n-2,\alpha}(x);
\]
and taking $x=0$ gives $\Delta_{\alpha}^2 Q_n(0) = c_2 \theta_2\theta_1 = c_2 \cc_2$, so
%\[
%c_2 = \frac{\Delta_{\alpha}^2Q_n(0)}{\cc_2}.
%\]
$c_2 = {\Delta_{\alpha}^2Q_n(0)}/{\cc_2}$.
Following this process, we get 
\[
\Delta_{\alpha}^n Q_n(x) = c_n \theta_n\theta_{n-1} \cdots \theta_1 f_{0,\alpha}(x),
\]
that evaluated at $x=0$ gives $\Delta_{\alpha}^n Q_n(0) = c_n\cc_n$, so $c_n = \Delta_{\alpha}^n Q_n(0)/\cc_n$.
\end{proof}

Actually, the Newton expansion \eqref{eq:Qnx} is very often written in the slightly more general way
\[
  Q_n(a+x) = \sum_{k=0}^n \frac{\Delta^kQ_n(a)}{k!}\, x^{\underline{k}}.
\]
In the Dunkl setting, this can be written as
\begin{equation}
\label{eq:aQnxD}
  \tau_x Q_n(a) = \tau_a Q_n(x) 
  = \sum_{k=0}^n \frac{\Delta_{\alpha}^kQ_n(a)}{\cc_{k,\alpha}}\, f_{k,\alpha}(x).
\end{equation}
But this expression is a simple consequence of~\eqref{eq:QnxD}. 
To show it, let us first note that, as a consequence of $\tau_a\tau_b = \tau_b\tau_a$ we have $\tau_a\Delta_\alpha = \tau_a\Delta_\alpha$, 
and $\tau_a\Delta_\alpha^k = \tau_a\Delta_\alpha^k$ for every~$k$.
Then, if we take the polynomial $R_n(x) = \tau_a Q_n(x)$, we have
$\Delta_\alpha^k R_n(x) = \Delta_\alpha^k \tau_a Q_n(x) = \tau_a \Delta_\alpha^k Q_n(x)$,
that evaluated at $x=0$ is
\[
\Delta_\alpha^k R_n(0) = \tau_a (\Delta_\alpha^k Q_n)(0) 
= \tau_0 (\Delta_\alpha^k Q_n)(a) = \Delta_\alpha^k Q_n(a).
\]
Finally, \eqref{eq:QnxD} applied to $R_n$ gives~\eqref{eq:aQnxD}.

%----------
\section{Bernoulli-Dunkl polynomials of the second kind}
\label{sec:BDsecond}
%----------

Recall that the Bernoulli polynomials of the second kind are defined by
\begin{equation*}
%\label{eq:fungenbk}
  \frac{t}{\log(1+t)}(1+t)^x = \sum_{k=0}^\infty b_k(x) \, \frac{t^k}{k!},
\end{equation*}
and that they satisfy the relation $b_k(x+1) - b_k(x) = k b_{k-1}(x)$. 
% forward differences
And, as explained in Section~\ref{sec:DAD}, instead of the forward differences, in the Dunkl context we must generalize classical cases with central differences. This already exists in the mathematical literature, see~\cite[\S\,6]{Tem}: the so-called central Bernoulli polynomials of the second kind $\{b_k^{\mathit{II}}(x)\}_{k=0}^{\infty}$ are defined by
\begin{equation}
\label{eq:bkII}
  \frac{t}{\log(t+\sqrt{1+t^2})} \big(t+\sqrt{1+t^2}\big)^x 
  = \sum_{k=0}^{\infty} b_k^{\mathit{II}}(x) \frac{t^k}{k!},
\end{equation}
and they satisfy 
\begin{equation}
\label{eq:bkII1-1}
  \frac{1}{2}(b_k^{\mathit{II}}(x+1) - b_k^{\mathit{II}}(x-1)) = k b_{k-1}^{\mathit{II}}(x).
\end{equation}

In our discrete Appell-Dunkl context, \eqref{eq:bkII1-1} is \eqref{eq:tau1-1pk} for $\alpha=-1/2$, so we must use \eqref{eq:defADdis} with a suitable $A(t)$ that extends the function $t/\log(t+\sqrt{1+t^2})$ in~\eqref{eq:bkII}. 
Since $\log(t+\sqrt{1+t^2}) = G_{\alpha}^{-1}(t)$ for $\alpha=-1/2$, we define the 
Bernoulli-Dunkl of the second kind as
\begin{equation}
\label{eq:bka}
  \frac{t}{G_{\alpha}^{-1}(t)} \Ea(x G_{\alpha}^{-1}(t)) 
  = \sum_{k=0}^{\infty} b_{k,\alpha}(x) \frac{t^k}{\cc_{k,\alpha}}, 
  \qquad x \in \RR.
\end{equation}
The first of these polynomials are
\begin{align*}
b_{0,\alpha}(x) &= 1, \qquad
b_{1,\alpha}(x) = x, \qquad
b_{2,\alpha}(x) = x^2 + \frac{\alpha+1}{\alpha+2}, \\
b_{3,\alpha}(x) &= x^3, \qquad
b_{4,\alpha}(x) = x^4 - 2 x^2
  - \frac{(\alpha+1)(3\alpha+10)}{(\alpha+2)(\alpha+3)}, \\
b_{5,\alpha}(x) &= x^5 - \frac{(4\alpha+12)x^3}{\alpha+2}, \\
b_{6,\alpha}(x) &= x^6
  - \frac{9(\alpha+3)x^4}{\alpha+2}
  + \frac{3(5\alpha+16) x^2}{\alpha+2}
  + \frac{(\alpha+1)(25\alpha^2+190\alpha+364)}{(\alpha+2)^2 (\alpha+4)}, \\
b_{7,\alpha}(x) &= x^7
  - \frac{12(\alpha+4)x^5}{\alpha+2}
  + \frac{6(\alpha+4)(6\alpha+19) x^3}{(\alpha+2)^2}.
\end{align*}

Let us now see how to extend a classical formula, namely
\[
  b_k(x) = x^{\underline{k}} - k! \sum_{j=0}^{k-1} \frac{b_j(x)}{j!} \frac{(-1)^{k-j}}{k+1-j},
  \qquad k \ge 1,
\]
a relation that can be used to compute the Bernoulli polynomials of the second kind in a recursive way,
and that, in particular, implies that
\[
  b_k(0) = - k! \sum_{j=0}^{k-1} \frac{(-1)^{k-j}b_j(0)}{(k+1-j)j!},
  \qquad k \ge 1.
\]
In the case of Bernoulli-Dunkl polynomials of the second kind, these recurrence relations are as follows:

\begin{theorem}
Let $\alpha>-1$ and $b_{k,\alpha}(x)$ be the Bernoulli-Dunkl polynomials of the second kind,
and denote $G_\alpha^{-1}(t) = \sum_{k=1}^{\infty} a_k t^k$ (the coefficients $a_k$ are identified in Theorem~\ref{teo:inver-Ga}).
Then, 
\begin{equation*}
%\label{eq:bk-recur}
b_{k,\alpha}(x) = f_{k,\alpha}(x) - \cc_{k,\alpha} \sum_{j=0}^{k-1} 
\frac{a_{k+1-j} b_{j,\alpha}(x)}{\cc_{j,\alpha}},
\quad k=1,2,\ldots,
\end{equation*}
and, for $x=0$,
\begin{equation*}
%\label{eq:bk-recur-0}
b_{k,\alpha}(0) = -\cc_{k,\alpha} \sum_{j=0}^{k-1} \frac{a_{k+1-j} b_{j,\alpha}(0)}{\cc_{j,\alpha}},
\quad k=1,2,\ldots.
\end{equation*}
\end{theorem}

\begin{proof}
It is enough to use Theorem~\ref{teo:fnfrompj} with $1/A(t) = G_\alpha^{-1}(t)$ and $G_\alpha^{-1}(0)=0$.
\end{proof}

Finally, in the same way than the generalized Bernoulli polynomials of the second kind and order $r$ are defined by
\begin{equation*}
  \left(\frac{t}{\log(1+t)}\right)^r (1+t)^x 
  = \sum_{k=0}^\infty b_k^{(r)}(x) \, \frac{t^k}{k!}
\end{equation*}
(see, for instance, \cite[\S\,2]{Car}), we define the generalized Bernoulli-Dunkl polynomials of the second kind and order $r$ as
\begin{equation}
\label{eq:bkar}
  \left(\frac{t}{G_{\alpha}^{-1}(t)}\right)^r \Ea(x G_{\alpha}^{-1}(t)) 
  = \sum_{k=0}^{\infty} b_{k,\alpha}^{(r)}(x) \frac{t^k}{\cc_{k,\alpha}}.
\end{equation}

One might expect that Bernoulli-Dunkl polynomials of the second kind 
(and order~$r$) would satisfy many identities and formulas corresponding to known 
properties of the classical Bernoulli polynomials of the second kind. 
Although the scope of this paper is not to look for these identities, 
we display here just a few of them to taste their flavor.

Let us start with a formula that relates the polynomials of order $r$ with
the polynomials of order $r-1$.

\begin{theorem}
\label{teo:Labka}
If we apply the Dunkl operator $\La$ to the Bernoulli-Dunkl polynomials of the second kind and order $r$, we have
\[
  \La b_{k,\alpha}^{(r)}(x) = \theta_{k,\alpha} b_{k-1,\alpha}^{(r-1)}(x),
  \quad k=1,2,\ldots;
\]
in particular, for $r=1$,
\[
  \La b_{k,\alpha} (x) = \theta_{k,\alpha} f_{k-1,\alpha}(x),
  \quad k=1,2,\ldots,
\]
where $f_{k,\alpha}(x)$ are the Dunkl factorial polynomials defined in~\eqref{eq:fallfka}.
\end{theorem}

\begin{proof}
Let us see the case $r=1$; to do it, let us apply $\Lax$ to both sides of \eqref{eq:bka}.
Then, by using \eqref{eq:LaEalx} we have
\[
  t \Ea(x G_{\alpha}^{-1}(t)) 
  = \sum_{k=0}^{\infty} (\La b_{k,\alpha})(x) \frac{t^k}{\cc_{k}}, 
  \qquad x \in \RR.
\]
Since $\La b_{0,\alpha} = 0$, we can write
\[
  \Ea(x G_{\alpha}^{-1}(t)) 
  = \sum_{k=1}^{\infty} (\La b_{k,\alpha})(x) \frac{t^{k-1}}{\cc_{k}} 
  = \sum_{k=0}^{\infty} (\La b_{k+1,\alpha})(x) \frac{t^{k}}{\cc_{k+1}}.
\]
Comparing to \eqref{eq:fallfka}, using $\theta_{k+1} = \cc_{k+1}/\cc_{k}$, and by the unicity of the analytic expansions,
we get $\La b_{k+1,\alpha} = \theta_{k+1} f_{k,\alpha}$.
The generalized case for arbitrary $r$ is completely similar but starting from~\eqref{eq:bkar}.
\end{proof}

The values at $0$ and $1$ of the classical Bernoulli polynomials 
of the second kind are related by means of
\[
  b_{k+1}(1) - b_{k+1}(0) = (k+1)b_k(0).
\]
In the Dunkl setting, also the value at $-1$ is involved in the corresponding relation.
Stated for the Bernoulli-Dunkl polynomials of the second kind, 
this property is as follows:

\begin{theorem}
\label{teo:rel01-1}
Let $\alpha>-1$ and $b_{k,\alpha}^{(r)}(x)$ be the Bernoulli-Dunkl polynomials of the second kind and order~$r$. Then, 
\begin{equation}
\label{eq:rel01-1}
  (\alpha+1) \big( b_{k,\alpha}^{(r)}(1)-b_{k,\alpha}^{(r)}(-1) \big)
  = \theta_{k,\alpha} b_{k-1,\alpha}^{(r)}(0),
  \quad k= 1,2,\ldots.
\end{equation}
\end{theorem}

\begin{proof}
Multiplying \eqref{eq:bkar} by $t$ and evaluating at $x=0$ we have
\begin{equation}
\label{eq:ev0}
\frac{t^{r+1}}{(G_{\alpha}^{-1}(t))^r}
= \sum_{k=1}^{\infty} \theta_k b_{k-1,\alpha}^{(r)}(0) \frac{t^k}{\cc_k}.
\end{equation}
Now, by evaluating $\Ea(xG_{\alpha}^{-1}(t))$ at $x=-1$ and $x=1$, 
$$
(\alpha+1) \left(\Ea(G_{\alpha}^{-1}(t))-\Ea(-G_{\alpha}^{-1}(t))\right)
= 2(\alpha+1) \sum_{k=0}^{\infty} \frac{(G_{\alpha}^{-1}(t))^{2k+1}}{\cc_{2k+1}}.
$$
But, by the construction of $G_{\alpha}^{-1}$ in Section~\ref{sec:DAD},
\[
(\alpha+1) \big(\Ea(G_{\alpha}^{-1}(t))-\Ea(-G_{\alpha}^{-1}(t))\big) = t,
\]
so the left hand side of \eqref{eq:ev0} becomes
\begin{align*}
\frac{t^{r+1}}{G_{\alpha}^{-1}(t)^r}
&= (\alpha+1) \bigg(\frac{t}{G_{\alpha}^{-1}(t)}\bigg)^r 
\big(\Ea(G_{\alpha}^{-1}(t))-\Ea(-G_{\alpha}^{-1}(t))\big) \\
&= (\alpha+1) \sum_{k=0}^{\infty} \frac{b_{k,\alpha}^{(r)}(1)-b_{k,\alpha}^{(r)}(-1)}{\cc_k} \,t^k.
\end{align*}
Equating coefficients with the right hand side of \eqref{eq:ev0} we obtain~\eqref{eq:rel01-1}.
\end{proof}

%----------
\begin{thebibliography}{99}
%----------

\bibitem{AdLe}
J. A. Adell and A. Lekuona,
Binomial convolution and transformations of Appell polynomials,
\textit{J. Math. Anal. Appl.}
\textbf{456} (2017), 16--33.

\bibitem{AlS}
W. A. Al-Salam,
Characterization theorems for orthogonal polynomials,
\textit{Orthogonal Polynomials: Theory and Practice}
(Columbus, OH, 1989, P.~Nevai, ed.), 1--24,
NATO Adv. Sci. Inst. Ser. C Math. Phys. Sci., 294,
Kluwer Acad. Publ., Dordrecht, 1990.
% NATO ASI Series, Volume 294 (1990)
% http://link.springer.com/book/10.1007/978-94-009-0501-6

\bibitem{AdJ} 
N. B. Andersen and M. de Jeu,
{Elementary proofs of Paley-Wiener theorems
for the Dunkl transform on the real line,}
\textit{Int. Math. Res. Not.}
\textbf{30} (2005), 1817--1831.

\bibitem{Ap}
T. M. Apostol,
On the Lerch zeta function,
\textit{Pacific J. Math.} \textbf{1} (1951), 161--167;
addendum, \textit{Pacific J. Math.} \textbf{2} (1952), 10.

\bibitem{Bab}
{J. Babini},
Polinomios generalizados de Bernoulli y sus correlativos,
%\textit{Revista Matem\'atica Hispano-Americana (2.\textsuperscript{a} serie)}
\textit{Rev. Mat. Hisp.-Amer. (4)}
\textbf{10} (1935), 23--25.

\bibitem{Car}
{L. Carlitz},
A note on Bernoulli and Euler polynomials of the second kind,
\textit{Scripta Math.}
\textbf{25} (1961), 323--330.

\bibitem{Cho} 
F. M. Cholewinski, 
\textit{The finite calculus associated with Bessel functions}, 
Contemporary Mathematics, 75, American Mathematical Society, Providence, RI, 1988.

\bibitem{CDPV}
\'O. Ciaurri, A. J. Dur\'an, M. P\'erez and J. L. Varona,
Bernoulli-Dunkl and Apostol-Euler-Dunkl polynomials with applications to series involving
zeros of Bessel functions,
\textit{J. Approx. Theory}
\textbf{235} (2018), 20--45.
% https://doi.org/10.1016/j.jat.2018.06.001

\bibitem{CMV}
\'O. Ciaurri, J. M\'inguez Ceniceros and J. L. Varona,
{Bernoulli-Dunkl and Euler-Dunkl polynomials and their generalizations},
\textit{Rev. R. Acad. Cienc. Exactas F\'is. Nat. Ser. A Mat. RACSAM} 
\textbf{113} (2019), 2853--2876.
% https://doi.org/10.1007/s13398-019-00662-z

\bibitem{CPRV}
\'O. Ciaurri, M. P\'erez, J. M. Reyes and J. L. Varona,
Mean convergence of Fourier-Dunkl series,
\textit{J. Math. Anal. Appl.}
\textbf{372} (2010), 470--485.
% http://www.unirioja.es/cu/jvarona/downloads/Fourier-Dunkl-series.pdf

%\bibitem{CV}
%\'O. Ciaurri and J. L. Varona,
%{A Whittaker-Shannon-Kotel'nikov sampling theorem related to the Dunkl transform},
%\textit{Proc. Amer. Math. Soc.}
%\textbf{135} (2007), 2939--2947.

\bibitem{CRV}
\'O. Ciaurri, L. Roncal and J. L. Varona,
The multiplier of the interval $[-1,1]$ for the Dunkl transform of arbitrary order on the real line, \textit{Integral Transforms Spec. Funct.} 
\textbf{26} (2015), 728--736.

\bibitem{DW}
% Dai, Feng (3-AB-MS); Wang, Heping (PRC-CAP-SM)
F. Day and H. Wang, 
Interpolation by weighted Paley-Wiener spaces associated with the Dunkl transform, 
\textit{J. Math. Anal. Appl.} \textbf{390} (2012), 556--572. 

\bibitem{Dil-NIST}
K. Dilcher,
Bernoulli and Euler polynomials,
\textit{NIST handbook of mathematical functions},
% (edited by F.~W. F.~Olver, D.~W. Lozier, R.~F. Boisvert and C.~W. Clark),
587--599, % Es el Chapter 24
U.S. Dept. Commerce, Washington, DC, 2010.
%% LONG CITE:
% National Institute of Standards and Technology, Washington, DC, 
% and Cambridge University Press, Cambridge, 2010.
Available online in \url{http://dlmf.nist.gov/24}

\bibitem{Du}
C. F. Dunkl,
{Differential-difference operators associated to reflection groups},
\textit{Trans. Amer. Math. Soc.}
\textbf{311} (1989), 167--183.

\bibitem{DPV}
A. J. Dur\'an, M. P\'erez and J. L. Varona,
Fourier-Dunkl system of the second kind and Euler-Dunkl polynomials,
\textit{J. Approx. Theory}
\textbf{245} (2019), 23--39.
% https://doi.org/10.1016/j.jat.2019.04.007

%\bibitem{EMOT2}
%A. Erd\'elyi, W. Magnus, F. Oberhettinger and F. G. Tricomi,
%\textit{Tables of integral transforms},
%vol.~II, McGraw-Hill, New York, 1954.

\bibitem{Gold}
L. Gold,
{Inverse Bessel functions: Solution for the zeros},
\textit{J. Math. and Phys.}
\textbf{36} (1957), 167--171.

\bibitem{GKP}
R. L. Graham, D. E. Knuth and O. Patashnik,
\textit{Concrete Mathematics: A Foundation for Computer Science}, 2nd ed.,
Addison-Wesley, Reading, MA, 1994.

\bibitem{Jeu}
M. F. E. de Jeu,
{The Dunkl transform},
\textit{Invent. Math.}
\textbf{113} (1993), 147--162.

\bibitem{Jor} % Charles Jordan
C. Jordan,
Sur les polynomes analogues aux polynomes de Bernoulli et sur des formules de sommation analogues \`a celle de MacLaurin-Euler,
%\textit{Acta Scientarium Mathematicarum} 
\textit{Acta Szeged} 
\textbf{4} (1929), 130--150.

\bibitem{Knuth}
D. E. Knuth, 
Two notes on notation, 
\textit{Amer. Math. Monthly} 
\textbf{99} (1992), 403--422. % no.~5
%arXiv:math/9205211, 
%doi:10.2307/2325085, 
%JSTOR 2325085 

\bibitem{MO}
H. Mejjaoli and A. Ould Ahmed Salem, 
On the Dunkl convolution equations in spaces of $\mathcal{D}_\alpha^p$ type, 
\textit{Integral Transforms Spec. Funct.}
\textbf{26} (2015), 619--634.

\bibitem{MSV}
D. Merlini, R. Sprugnoli and M. C. Verri,
The Cauchy numbers,
\textit{Discrete Math.}
\textbf{306} (2006), 1906--1920.

% MR0043339 (realmente es un ``reprinted from the first edition of 1933'')
\bibitem{MT}
L. M. Milne-Thomson, 
\textit{The Calculus of Finite Differences}, 
Macmillan and Co., London, 1951.

\bibitem{NRV-Brno}
L. M. Navas, F. J. Ruiz and J. L. Varona, 
Existence and reduction of generalized Apostol-Bernoulli, Apostol-Euler and Apostol-Genocchi polynomials, 
\textit{Arch. Math. (Brno)} 
\textbf{52} (2019), 157--165. 

\bibitem{Nor}
N. E. N{\o}rlund, 
M\'emoire sur les polyn\^omes de Bernoulli, 
\textit{Acta. Math.} 
\textbf{43} (1922), 121--196.

%% Aqu\'i define Euler generalizado:
%N. E. N{\o}rlund, 
%Vorlesungen \"uber Differenzenrechnung, 1st ed., 
%Springer-Verlag, Berlin-Heidelberg, 1924.

\bibitem{NS}
A. Nowak and K. Stempak, 
Relating transplantation and multipliers for Dunkl and Hankel transforms, 
\textit{Math. Nachr.} 
\textbf{281} (2008), 1604--1611. 

\bibitem{OlMa-NIST}
F. W. J. Olver and L. C. Maximon,
Bessel Functions,
\textit{NIST handbook of mathematical functions},
% (edited by F.~W. F.~Olver, D.~W. Lozier, R.~F. Boisvert and C.~W. Clark),
215--286, % Es el Chapter 10
U.S. Dept. Commerce, Washington, DC, 2010.
%% LONG CITE:
% National Institute of Standards and Technology, Washington, DC, 
% and Cambridge University Press, Cambridge, 2010.
Available online in \url{http://dlmf.nist.gov/10}

\bibitem{Ph} % Lo cita OEIS
G. M. Phillips, 
Gregory's method for numerical integration, 
\textit{Amer. Math. Monthly}
\textbf{79} (1972), 270--274.

\bibitem{RP-BSMA}
{J. Rey Pastor},
Polinomios correlativos de los de Bernoulli,
%\textit{Bolet\'in del Seminario Matem\'atico Argentino}
\textit{Bolet\'in Seminario Mat. Argentino}
\textbf{1} (1929), no.~3, 1--10.

\bibitem{Ros}
M. Rosenblum,
{Generalized Hermite polynomials and the Bose-like oscillator calculus},
\textit{Oper. Theory Adv. Appl.}
\textbf{73} (1994), 369--396.

\bibitem{Rosler-LN}
M. R\"osler,
Dunkl operators: theory and applications,
\textit{Orthogonal polynomials and special functions 
(Leuven, 2002)}, 93--135, 
Lecture Notes in Math., 1817, Springer, Berlin, 2003.

\bibitem{Rosler-TAMS}
M. R\"osler,
A positive radial product formula for the Dunkl kernel,
\textit{Trans. Amer. Math. Soc.}
\textbf{355} (2003), 2413--2438.

\bibitem{Tem}
P. Tempesta,
On Appell sequences of polynomials of Bernoulli and Euler type,
\textit{J. Math. Anal. Appl.} 
\textbf{341} (2008), 1295--1310.

\bibitem{ThX}
S. Thangavelu and Y. Xu,
Riesz transform and Riesz potentials for Dunkl transform,
\textit{J. Comput. Appl. Math.}
\textbf{199} (2007), 181--195.

\bibitem{Tr}
K. Trim\`eche,
Paley-Wiener theorems for the Dunkl transform and Dunkl translation operators.
\textit{Integral Transform Spec. Funct.} 
\textbf{13} (2002), 17--38.

\bibitem{Wat}
G. N. Watson,
\textit{A treatise on the theory of Bessel functions},
Cambridge University Press, Cambridge, 1958.

\bibitem{WuPan}
% Ming Wu and Hao Pan, 
M. Wu and H. Pan, 
Sums of products of Bernoulli numbers of the second kind, 
\textit{Fib. Quart.} \textbf{45} (2007), 146--150.

\bibitem{Zh}
F.-Z. Zhao, % Feng-Zhen Zhao
Sums of products of Cauchy numbers,
\textit{Discrete Math.}
\textbf{309} (2009), 3830--3842.

%----------
\end{thebibliography}
%----------

%----------
\end{document}
%---------- 
